\documentclass[11pt]{article}

\usepackage[margin=1in]{geometry}
\usepackage{amsmath,amssymb}
\usepackage{booktabs}
\usepackage{siunitx}
\usepackage{xcolor}
\usepackage{hyperref}
\usepackage{enumitem}
\usepackage{caption}

\sisetup{round-mode=places, round-precision=4, table-format=-1.4}
\hypersetup{colorlinks=true, linkcolor=blue!50!black, urlcolor=blue!50!black}
\setlist{itemsep=2pt, topsep=4pt}
\usepackage{microtype}
\newcommand{\codeus}{\textunderscore\allowbreak}
\newcommand{\code}[1]{{\ttfamily\let\_\codeus #1}}

\title{Kron-Reduced vs.\ Full Graph Laplacians for\\
       Water-Quality Regression on the Colorado River Network}
\author{Andrew Chincea \\ TDA Lab}
\date{November 2026}

\begin{document}
\maketitle

\begin{abstract}
The graph-Laplacian water-quality models for the Colorado basin use a Kron (Schur-complement)
reduction of the 67{,}401-node river network onto the 26 nodes with water-quality (WQ) stations.
We test whether this reduction changes the model. For five MERIT-based edge weightings, on both
the uncompressed and the compressed (67-node) graph, the Kron and full Laplacians agree to within
$10^{-8}$ on every quantity the reduction should preserve. Re-fitting the nested
leave-one-node-out (LOO) regression with the full Laplacian changes $R^2$ by at most
$10^{-10}$ and selects identical hyperparameters in every fold. Kron reduction is therefore exact,
and about $150\times$ faster. We also correct several problems in the original analysis. The most
consequential is that feature sets chosen by looking at outer-LOO results overstated performance:
by about $0.13$ in $R^2$ for specific conductance and $0.02$--$0.08$ for DO, pH and turbidity.
With selection done inside the inner loop, the graph term clearly helps for pH and specific
conductance, while seasonal and MERIT predictors matter much more for water temperature, DO and
turbidity.
\end{abstract}

% =====================================================================
\section{Setup}

\paragraph{Data.} Hourly USGS measurements at 28 stations, mapped to 26 river-network nodes.
Depending on the target, 17--24 nodes have usable data
(DO 17, water temperature 24, pH 19, specific conductance 24, turbidity 17).
The network is the largest connected component of the MERIT-derived Colorado graph:
67{,}401 nodes and 67{,}400 edges, each edge about 70--120\,m long.

\paragraph{Model.} For node $i$ and hour $t$, the prediction is
$\hat y_{it} = \beta_0 + z_{it}^\top\beta + g_i$, where $z_{it}$ are standardized predictors and
$g_i$ is a node offset. The fit minimizes
\begin{equation}
  \sum_{i\in\text{train}} \frac{1}{n_i}\sum_t \bigl(y_{it}-\beta_0-z_{it}^\top\beta-g_i\bigr)^2
  \;+\; \lambda_\beta \sum_{k\ge 1}\beta_k^2 \;+\; \lambda\, g^\top \tilde L\, g,
  \qquad \text{s.t. } \textstyle\sum_{i\in \mathrm{WQ}} g_i = 0 ,
  \label{eq:model}
\end{equation}
with $g^\top L g = \sum_{(i,j)} w_{ij}(g_i-g_j)^2$ and $\tilde L = L/s$. The scale $s$ is the mean
positive diagonal of the Kron matrix. The weight $1/n_i$ makes each station count equally.
A held-out node has no data term, so its offset is set by the smoothness penalty alone.
The hyperparameters $\lambda$ and $\lambda_\beta$ (and, where stated, the edge weighting) are tuned
by an inner LOO. Performance is the node-weighted $R^2$ of the outer-LOO predictions.

\paragraph{Edge weightings.} $w_{ij} = \exp(\text{score}_{ij})/d_{ij}$, where
$d_{ij}$ is the edge length in km and the score mixes z-scored log upstream area, log width and
log slope. There are five schemes: \code{distance\_only}, \code{area}, \code{width},
\code{slope} and \code{area\_width}.

\paragraph{Kron reduction.} Partitioning nodes into WQ nodes $B$ and the rest $I$,
\begin{equation}
  L_{\mathrm{kron}} = L_{BB} - L_{BI} L_{II}^{-1} L_{IB} \in \mathbb{R}^{26\times 26}.
\end{equation}
The non-WQ nodes have no data term in \eqref{eq:model}, so minimizing over their $g$ in closed form
leaves exactly $g_B^\top L_{\mathrm{kron}} g_B$. The two formulations should therefore be the same
model. This report checks that numerically.

% =====================================================================
\section{Are the Laplacians the same?}

The Kron and full matrices have different sizes, so we compare what Kron reduction must preserve:
\begin{itemize}
  \item \textbf{Effective resistance} between every pair of WQ nodes, computed on the full
        Laplacian by sparse solves and on $L_{\mathrm{kron}}$.
  \item \textbf{Dirichlet energy:} for random WQ values $x_B$ with harmonic extension $x$,
        $x^\top L x$ should equal $x_B^\top L_{\mathrm{kron}} x_B$.
  \item \textbf{Compression:} Kron(uncompressed) vs Kron(compressed), compared directly.
\end{itemize}
As a reference for what a \emph{wrong} reduction looks like, we also compare against the
sub-matrix $L_{BB}$, i.e.\ simply deleting the non-WQ nodes.

\begin{table}[h]
\centering
\caption{Largest relative errors, uncompressed graph. All values are at rounding-error level,
except the deliberately wrong $L_{BB}$ reference.}
\label{tab:matrix}
\begin{tabular}{lcccc}
\toprule
Weighting & Eff.\ resistance & Energy & Kron unc.\ vs comp. & Naive $L_{BB}$ (reference) \\
\midrule
\code{distance\_only} & $1.9\times10^{-9}$ & $2.4\times10^{-12}$ & $6.7\times10^{-13}$ & 205 \\
\code{area}           & $4.7\times10^{-9}$ & $1.5\times10^{-12}$ & $5.1\times10^{-13}$ & 368 \\
\code{width}          & $4.1\times10^{-10}$& $3.1\times10^{-12}$ & $6.4\times10^{-13}$ & 253 \\
\code{slope}          & $7.3\times10^{-9}$ & $8.8\times10^{-12}$ & $2.6\times10^{-12}$ & 257 \\
\code{area\_width}    & $6.0\times10^{-9}$ & $2.6\times10^{-12}$ & $8.1\times10^{-13}$ & 359 \\
\bottomrule
\end{tabular}
\end{table}

On the compressed graph, the resistance and energy errors are below $1.5\times10^{-12}$ for every
weighting. Compression is also exact. Each compressed edge replaces a chain of resistors in series,
and series resistances add, so the total resistance between stations is unchanged.

% =====================================================================
\section{Does Kron change the regression?}

The nested LOO was re-run with four Laplacians, using exactly the same solver code:
\begin{center}
\begin{tabular}{lll}
\toprule
Variant & Laplacian in \eqref{eq:model} & Size \\
\midrule
\code{kron\_uncompressed} & Kron of the 67k graph (original method) & $26\times26$ \\
\code{kron\_compressed}   & Kron of the 67-node graph & $26\times26$ \\
\code{full\_uncompressed} & full Laplacian, $g$ on every node & $67{,}401\times67{,}401$ \\
\code{full\_compressed}   & full Laplacian, $g$ on every node & $67\times67$ \\
\bottomrule
\end{tabular}
\end{center}
The full variants never form a Schur complement. They solve the complete $(\beta, g)$ system by
block elimination with a sparse factorization. Each full Laplacian is divided by the same $s$ as
its Kron reduction, so $\lambda$ has the same meaning in all four.

\begin{table}[h]
\centering
\caption{Outer-LOO $R^2$, chemistry predictors only. ``Graph, $1/d$'' is \code{distance\_only}. The last column is the largest change in
$R^2$ between \code{kron\_uncompressed} and any other variant, over both graph modes.}
\label{tab:r2}
\begin{tabular}{lS S S c}
\toprule
Target & {Linear} & {Graph, $1/d$} & {Graph, tuned weighting} & {Max $|\Delta R^2|$} \\
\midrule
DO concentration     &  0.6208 &  0.6135 &  0.6065 & $3.4\times10^{-11}$ \\
Water temperature    & -0.0556 & -0.0346 & -0.0575 & $9.9\times10^{-11}$ \\
pH                   & -0.0352 &  0.0742 &  0.0742 & $1.6\times10^{-11}$ \\
Specific conductance & -0.0891 &  0.0907 &  0.1062 & $2.2\times10^{-11}$ \\
Turbidity            &  0.0228 & -0.0486 & -0.0360 & $4.4\times10^{-11}$ \\
\bottomrule
\end{tabular}
\end{table}

\begin{itemize}
  \item \textbf{Identical models.} All four variants chose the same (weighting, $\lambda$, $\lambda_\beta$)
        in 100\% of outer folds for every target. $R^2$ differs by at most $10^{-10}$.
  \item \textbf{Reproduction.} \code{kron\_uncompressed} reproduces the original notebook's saved
        results (\code{full\_graph\_nested\_loo\_results.csv} and the \code{Linear + Laplacian}
        rows of \code{all\_models\_node\_loo\_results.csv}) to within $1.3\times10^{-12}$.
  \item \textbf{Cost.} Kron took $0.1$\,s per outer fold vs $15.2$\,s for the full 67k
        Laplacian (26 minutes in total), about $150\times$ faster.
\end{itemize}

\paragraph{When the graph term helps.} With chemistry predictors only, the graph term improves pH
($+0.11$) and specific conductance ($+0.18$ to $+0.20$). Both are plausibly controlled by geology
and therefore spatially structured. It is neutral or slightly harmful for DO, water temperature and
turbidity. Tuning the MERIT weighting instead of using $1/d$ changes $R^2$ by at most $0.02$. The
weighting chosen is also unstable from fold to fold for DO (\code{width} in 13 of 17 folds) and for
water temperature (\code{area} 15, \code{slope} 5, \code{area\_width} 4 of 24). For pH,
\code{distance\_only} was chosen in all 19 folds.

% =====================================================================
\section{Corrections to the original analysis}

\begin{table}[h]
\centering
\caption{Problems found in \code{largest\_uncompressed\_laplacian.ipynb} and how they were handled.}
\begin{tabular}{p{0.47\textwidth}p{0.45\textwidth}}
\toprule
Problem & Fix \\
\midrule
pH-transform section used the \emph{unscaled} Kron Laplacian ($\lambda$ about $21\times$ weaker) &
scaled Laplacian; unscaled run alongside for comparison \\
Static feature sets and Fourier $K$ chosen from outer-LOO results, then re-scored on the same folds &
selection moved inside the inner LOO \\
``Discharge'' edge feature is a linear function of $\log$ UPA, so it duplicates \code{area};
its fit also used held-out stations &
discharge removed \\
HAND is $0$ at all 26 WQ nodes (river pixels) & HAND removed from the predictors \\
Comparison tables used hard-coded $R^2$ values from earlier runs & all tables computed from the run \\
\bottomrule
\end{tabular}
\end{table}

\subsection{pH transforms with the scaled Laplacian}

Predictors are specific conductance plus $K=1$ Fourier terms.

\begin{table}[h]
\centering
\caption{pH nested-LOO $R^2$. The scaling fix changes graph $R^2$ by only $-0.002$, because tuning
$\lambda$ over a wide grid absorbs the constant factor.}
\begin{tabular}{lS S S S}
\toprule
Target transform & {Linear} & {Graph, scaled (fixed)} & {Graph, unscaled (original)} & {Change} \\
\midrule
raw pH          &  0.0163 & 0.1300 & 0.1315 & -0.0016 \\
z-scored pH     &  0.0163 & 0.1300 & 0.1315 & -0.0016 \\
$H^+=10^{7-\mathrm{pH}}$ & -0.0662 & 0.0930 & 0.0954 & -0.0024 \\
\bottomrule
\end{tabular}
\end{table}

Raw and z-scored pH give identical results, as expected for an affine transform. Fitting $H^+$
is worse, and no $H^+$ predictions had to be clipped. Raw pH should be kept.

\subsection{Feature selection inside the nested LOO}

In each outer fold, the feature set was chosen by inner-LOO MSE alone, separately for the linear
and graph models (graph: Kron Laplacians, with weighting, $\lambda$ and $\lambda_\beta$ tuned). The
30 candidates are chemistry plus Fourier $K\in\{1,\dots,5\}$, combined with one of six static
options: none, one of elevation, $\log$ UPA, $\log$ width or local slope, or all four. The same run
gives each candidate's own nested-LOO $R^2$. Taking the best of those reproduces what the original
analysis effectively reported, so the difference measures how much it was overstated.

\begin{table}[h]
\centering
\caption{Outer-LOO $R^2$ with seasonal and MERIT predictors. ``Biased'' is the best candidate chosen
after seeing outer results. ``Fair'' is the feature set chosen inside the inner loop, and is the
number to report. ``Baseline'' is $K=1$ with no static features.}
\label{tab:selection}
\begin{tabular}{llS S S S}
\toprule
Target & Model & {Baseline} & {Biased} & {\textbf{Fair}} & {Overstated} \\
\midrule
Water temperature    & linear &  0.6781 &  0.8599 &  0.8573 & 0.0026 \\
                     & graph  &  0.7394 &  0.8591 &  0.8568 & 0.0024 \\
DO concentration     & linear &  0.6908 &  0.7260 &  0.7109 & 0.0151 \\
                     & graph  &  0.6831 &  0.7220 &  0.7009 & 0.0211 \\
pH                   & linear &  0.0163 &  0.0268 & -0.0553 & 0.0821 \\
                     & graph  &  0.1300 &  0.1439 &  0.0977 & 0.0461 \\
Specific conductance & linear & -0.0793 & -0.0564 & -0.1690 & 0.1126 \\
                     & graph  &  0.1462 &  0.2674 &  0.1403 & 0.1271 \\
Turbidity            & linear &  0.0554 &  0.1954 &  0.1443 & 0.0511 \\
                     & graph  & -0.0084 &  0.1064 &  0.0740 & 0.0324 \\
\bottomrule
\end{tabular}
\end{table}

\begin{itemize}
  \item \textbf{Water temperature.} A large and genuine gain, from $0.68$ to $0.86$, with almost
        no overstatement. All four MERIT features were chosen in every fold. With them included,
        the graph term adds nothing, because the node-level features already explain the offsets
        the Laplacian was capturing.
  \item \textbf{DO.} A modest genuine gain, to $0.70$--$0.71$, driven by $K=2$--$4$ seasonal terms
        with $\log$ UPA (graph, 13 of 17 folds) or elevation (linear, 14 of 17 folds).
        The graph is slightly worse than linear.
  \item \textbf{Specific conductance.} The largest overstatement. Choosing local slope after seeing
        outer results inflated graph $R^2$ from a fair $0.14$ to $0.27$, and the fair value is no
        better than the plain $K=1$ baseline ($0.146$). The graph term is still the only model with
        positive $R^2$. Local slope was chosen in 22 of 24 folds, so it may be real, but it does
        not improve held-out accuracy.
  \item \textbf{pH.} The graph term helps ($0.10$ vs $-0.06$ for linear), but $R^2$ stays low and
        the chosen features vary from fold to fold.
  \item \textbf{Turbidity.} Linear with $K=5$ seasonal terms and no static features is best
        ($0.14$). The graph term hurts.
\end{itemize}

% =====================================================================
\section{Conclusions}

\begin{enumerate}
  \item \textbf{Kron reduction is exact for this model.} It matches the full 67{,}401-node
        Laplacian to rounding error for every weighting, and gives identical hyperparameters and
        $R^2$. It is about $150\times$ faster, so the $26\times26$ Kron matrix should remain
        the default. Compressing the graph changes nothing either.
  \item \textbf{The value of the graph term depends on the target.} It clearly helps for pH and,
        to a lesser extent, specific conductance. For water temperature, DO and turbidity,
        seasonal and MERIT predictors matter far more than the Laplacian.
  \item \textbf{Report the fair numbers.} The fair $R^2$ values in Table~\ref{tab:selection}
        should replace earlier figures. In particular, specific conductance should be reported at
        about $0.14$, not $0.22$--$0.27$.
\end{enumerate}

\paragraph{Caveats.} With only 17--24 stations, differences in $R^2$ below roughly $0.03$--$0.05$
are within noise, and no standard errors were computed. For pH and specific conductance the chosen
features vary between folds, another sign that those gains are fragile. The daily$\times$annual
interaction terms and separate daily and annual $K$ from the original water-temperature section
were not among the candidates. In the original they changed predictions by at most $0.009$\,\textdegree C.

\paragraph{Reproducibility.} All results come from \code{largest\_uncompressed\_laplacian\_fixed.ipynb}
and are saved in \code{results/laplacian\_fixed/}. The standalone check script is
\code{compare\_kron\_vs\_full\_laplacian.py}.

\end{document}
